\section{Randomized --- Fingerprint}
%%% generated by the blueprint pipeline from TCSlib.Complexity.Randomized.Fingerprint
\subsection*{Randomized --- Fingerprint: Introduction}

Declarations of the Lean module \texttt{TCSlib.Complexity.Randomized.Fingerprint}.

\subsection*{Randomized --- Fingerprint: Results}

\begin{theorem}[Prime-counting exponential bound]
\lean{Randomized.four_pow_le_pow_primeCounting}
\label{Randomized.four_pow_le_pow_primeCounting}
\leanok
\difficulty{5}
For every positive natural number $n$, $4^n\le (2n)^{\pi(2n)+1}$, where $\pi(x)$ denotes
the number of primes at most $x$.
\end{theorem}

\begin{theorem}[Prime counting lower bound]
\lean{Randomized.two_pow_le_mul_primeCounting}
\label{Randomized.two_pow_le_mul_primeCounting}
\leanok
\uses{Randomized.four_pow_le_pow_primeCounting}
\difficulty{3}
For every natural number $L\ge1$, the number of primes at most $2^L$, plus one,
multiplied by $L$ is at least $2^L$.
\end{theorem}

\begin{theorem}[Prime factors bounded by the integer]
\lean{Randomized.two_pow_card_primeFactors_le}
\label{Randomized.two_pow_card_primeFactors_le}
\leanok
\difficulty{3}
For every nonzero natural number $k$, the number of distinct prime divisors of $k$ is at
most $\log_2 k$.
\end{theorem}

\begin{theorem}[Prime Count Bounded by Nondivisors]
\lean{Randomized.primeCounting_le_card_not_dvd}
\label{Randomized.primeCounting_le_card_not_dvd}
\leanok
\uses{Randomized.two_pow_card_primeFactors_le}
\difficulty{5}
Let $y\ne0$ be an integer and let $B,N$ be natural numbers such that $|y|\le2^B$. The
number of primes at most $N$ is no greater than the number of integers $v$ with $0\le
v<N$ for which $v+1$ does not divide $y$, plus $B$.
\end{theorem}

\begin{theorem}[An exponential inequality]
\lean{Randomized.four_mul_le_four_pow}
\label{Randomized.four_mul_le_four_pow}
\leanok
\difficulty{4}
For every natural number $M\ge 12$, $4M(1+5M2^M)\le 2^{2M}$.
\end{theorem}

\begin{theorem}[Nondivisibility fingerprint probability]
\lean{Randomized.fingerprint}
\label{Randomized.fingerprint}
\leanok
\uses{Complexity.LogProg.bitsVal, Randomized.four_mul_le_four_pow, Randomized.primeCounting_le_card_not_dvd, Randomized.randProb, Randomized.randProb_bitsVal, Randomized.two_pow_le_mul_primeCounting}
\difficulty{5}
Let $M\ge12$ be an integer and let $y$ be a nonzero integer with $|y|\le 2^{5M2^M}$.
Choose a Boolean string $r$ uniformly from strings of length $2M$, and interpret it as
the nonnegative integer whose bits, read from least significant to most significant,
have place values $1,2,4,\ldots$. The probability that this integer plus one does not
divide $y$ is at least $1/(4M)$.
\end{theorem}

\begin{theorem}[High Probability of a Nondividing Block Value]
\lean{Randomized.prob_exists_not_dvd}
\label{Randomized.prob_exists_not_dvd}
\leanok
\uses{Complexity.LogProg.bitsVal, Randomized.blockVal, Randomized.fingerprint, Randomized.one_sub_pow_le_half_pow, Randomized.randProb, Randomized.randProb_forall_blocks, Randomized.randProb_nonneg, Randomized.randProb_not}
\difficulty{5}
Let $M\ge12$ be a natural number and let $y\ne0$ be an integer with
$|y|\le2^{5M\cdot2^M}$. Choose uniformly a bit string $s$ of length $16M^2$. For each
$j=0,\ldots,8M-1$, let $r_j$ be the integer represented in little-endian order by the
block of $2M$ bits of $s$ starting at position $2Mj$. Then with probability at least
$3/4$, there is some such $j$ for which $r_j+1$ does not divide $y$.
\end{theorem}

\begin{definition}[Input bit count]
\lean{Randomized.inputBits}
\label{Randomized.inputBits}
\leanok
For a natural number $m$, the input bit count is $m+4$.
\end{definition}

\begin{definition}[Modified parameter]
\lean{Randomized.modParam}
\label{Randomized.modParam}
\leanok
For a natural number $m$, the modified parameter is $m+12$.
\end{definition}

\begin{definition}[Number of trials]
\lean{Randomized.numTrials}
\label{Randomized.numTrials}
\leanok
\uses{Randomized.modParam}
For a natural number $m$, the number of trials is eight times the modified parameter
associated with $m$.
\end{definition}

\begin{definition}[Test length]
\lean{Randomized.testLen}
\label{Randomized.testLen}
\leanok
\uses{Randomized.inputBits, Randomized.modParam, Randomized.numTrials}
For natural numbers $n$ and $m$, the test length is $n$ times the input bit count for
$m$, plus the number of trials for $m$ times twice the modified parameter for $m$.
\end{definition}

\begin{definition}[Circuit test point]
\lean{Randomized.AlgCircuit.testPoint}
\statementsource{arora-barak-ch07-randomized-computation}{pdf-26d95080dbde-p005-b003}
\label{Randomized.AlgCircuit.testPoint}
\leanok
\uses{Randomized.AlgCircuit, Randomized.AlgCircuit.size, Randomized.blockVal, Randomized.inputBits}
Let $C$ be an arithmetic circuit on $n$ inputs and let $r$ be a finite Boolean string.
The test point associated with $C$ and $r$ is the point in $\mathbb N^n$ whose $i$th
coordinate is one plus the value of the block of $r$ assigned to input $i$, using blocks
of length determined by the size of $C$.
\end{definition}

\begin{definition}[Circuit test modulus]
\lean{Randomized.AlgCircuit.testModulus}
\statementsource{arora-barak-ch07-randomized-computation}{pdf-26d95080dbde-p005-b004}
\label{Randomized.AlgCircuit.testModulus}
\leanok
\uses{Randomized.AlgCircuit, Randomized.AlgCircuit.size, Randomized.blockVal, Randomized.modParam}
Let $C$ be an arithmetic circuit, let $s$ be a finite Boolean string, and let $j$ be a
natural number. The test modulus at $j$ is one plus the value of the $j$th block of $s$,
with blocks of length twice the modulus parameter determined by the size of $C$.
\end{definition}

\begin{definition}[Core modular zero test]
\lean{Randomized.AlgCircuit.ZeroTestCore}
\statementsource{arora-barak-ch07-randomized-computation}{pdf-26d95080dbde-p005-b003}
\label{Randomized.AlgCircuit.ZeroTestCore}
\leanok
\uses{Randomized.AlgCircuit, Randomized.AlgCircuit.eval, Randomized.AlgCircuit.size, Randomized.AlgCircuit.testModulus, Randomized.AlgCircuit.testPoint, Randomized.inputBits, Randomized.numTrials}
Let $C$ be an arithmetic circuit on $n$ inputs and let $r$ be a finite Boolean string.
For every trial index $j$ among the number of trials determined by the size of $C$,
evaluate $C$ modulo the test modulus obtained from the portion of $r$ after its first
$n$ input blocks, at the test point obtained from those first blocks. The value must be
zero in the corresponding residue class ring.
\end{definition}

\begin{definition}[Circuit zero test]
\lean{Randomized.AlgCircuit.zeroTest}
\statementsource{Vad12}{
  pdf-bf2ce5422088-p016-b003,
  pdf-bf2ce5422088-p044-b006
}
\statementsource{arora-barak-ch07-randomized-computation}{
  pdf-26d95080dbde-p005-b003,
  pdf-26d95080dbde-p005-b004
}
\label{Randomized.AlgCircuit.zeroTest}
\leanok
\uses{Randomized.AlgCircuit, Randomized.AlgCircuit.ZeroTestCore, Randomized.AlgCircuit.size, Randomized.testLen}
For a natural number $n$, an algebraic circuit $C$ with $n$ input variables, and a
Boolean list $r$, define the result to be whether the circuit's zero-test core returns
true on the prefix of $r$ whose length is the prescribed test length for $n$ and the
size of $C$.
\end{definition}

\begin{theorem}[Evaluation modulo a natural number]
\lean{Randomized.AlgCircuit.eval_natCast_zmod}
\label{Randomized.AlgCircuit.eval_natCast_zmod}
\leanok
\uses{Randomized.AlgCircuit, Randomized.AlgCircuit.eval, Randomized.AlgCircuit.map_eval}
\difficulty{3}
Let $C$ be an algebraic circuit on $n$ inputs, let $a$ assign a natural number to each
input, and let $k$ be a natural number. Evaluating $C$ on the reductions of the inputs
modulo $k$ gives the reduction modulo $k$ of its integer evaluation on $a$.
\end{theorem}

\begin{theorem}[Zero test on the zero polynomial]
\lean{Randomized.AlgCircuit.zeroTest_of_toMvPolynomial_eq_zero}
\label{Randomized.AlgCircuit.zeroTest_of_toMvPolynomial_eq_zero}
\leanok
\uses{Randomized.AlgCircuit, Randomized.AlgCircuit.aeval_toMvPolynomial, Randomized.AlgCircuit.toMvPolynomial, Randomized.AlgCircuit.zeroTest}
\difficulty{3}
Let $C$ be an algebraic circuit on $n$ inputs whose associated polynomial over
$\mathbb{Z}$ is zero. Then for every Boolean string $r$, the circuit's zero test on $r$
returns true.
\end{theorem}

\begin{theorem}[Prob eval ne zero]
\lean{Randomized.AlgCircuit.prob_eval_ne_zero}
\statementsource{Vad12}{
  pdf-bf2ce5422088-p017-b001,
  pdf-bf2ce5422088-p017-b002
}
\statementsource{arora-barak-ch07-randomized-computation}{pdf-26d95080dbde-p005-b002}
\label{Randomized.AlgCircuit.prob_eval_ne_zero}
\leanok
\uses{Randomized.AlgCircuit, Randomized.AlgCircuit.eval, Randomized.AlgCircuit.eval_toMvPolynomial, Randomized.AlgCircuit.size, Randomized.AlgCircuit.testPoint, Randomized.AlgCircuit.toMvPolynomial, Randomized.AlgCircuit.totalDegree_toMvPolynomial_le, Randomized.inputBits, Randomized.randProb, Randomized.randProb_blocks, Randomized.schwartz_zippel}
\difficulty{6}
% dataset-pipeline: stub -- prose pending build_informal_statements.py
Statement pending regeneration.
\end{theorem}

\begin{theorem}[Circuit evaluation bound at test points]
\lean{Randomized.AlgCircuit.natAbs_eval_testPoint_le}
\label{Randomized.AlgCircuit.natAbs_eval_testPoint_le}
\leanok
\uses{Randomized.AlgCircuit, Randomized.AlgCircuit.abs_eval_le_pow, Randomized.AlgCircuit.eval, Randomized.AlgCircuit.size, Randomized.AlgCircuit.testPoint, Randomized.blockVal_lt, Randomized.inputBits, Randomized.modParam}
\difficulty{4}
Let $C$ be a finite arithmetic circuit on $n$ inputs, and let $u$ be any finite Boolean
string. Form an input whose $i$th coordinate is one plus the value, in
least-significant-bit-first convention, of the block of $|C|+4$ bits of $u$ starting at
position $i(|C|+4)$; missing bits contribute nothing. Then the absolute value of the
integer output of $C$ at this input is at most $2^{5(|C|+12)2^{|C|+12}}$, where $|C|$
denotes the number of gates.
\end{theorem}

\begin{theorem}[Soundness probability of the zero test]
\lean{Randomized.AlgCircuit.zeroTest_sound}
\label{Randomized.AlgCircuit.zeroTest_sound}
\leanok
\uses{Randomized.AlgCircuit, Randomized.AlgCircuit.ZeroTestCore, Randomized.AlgCircuit.eval, Randomized.AlgCircuit.eval_natCast_zmod, Randomized.AlgCircuit.natAbs_eval_testPoint_le, Randomized.AlgCircuit.prob_eval_ne_zero, Randomized.AlgCircuit.size, Randomized.AlgCircuit.testModulus, Randomized.AlgCircuit.testPoint, Randomized.AlgCircuit.toMvPolynomial, Randomized.AlgCircuit.zeroTest, Randomized.blockVal, Randomized.inputBits, Randomized.modParam, Randomized.numTrials, Randomized.prob_exists_not_dvd, Randomized.randProb, Randomized.randProb_append_ge, Randomized.randProb_congr, Randomized.randProb_take, Randomized.testLen}
\difficulty{5}
Let $C$ be an algebraic circuit on $n$ inputs whose associated polynomial over
$\mathbb{Z}$ is nonzero, and let $N$ be a nonnegative integer at least the test length
for $C$. When an $N$-bit string is chosen uniformly at random, the probability that the
zero test returns false is at least $2/3$.
\end{theorem}

\begin{definition}[Zero-polynomial verifier]
\lean{Randomized.zeropVerifier}
\statementsource{arora-barak-ch07-randomized-computation}{
  pdf-26d95080dbde-p005-b003,
  pdf-26d95080dbde-p005-b004
}
\label{Randomized.zeropVerifier}
\leanok
\uses{Randomized.AlgCircuit.WellFormed, Randomized.AlgCircuit.decode, Randomized.AlgCircuit.zeroTest}
Given Boolean lists $x$ and $r$, decode $x$ as an algebraic circuit encoding. Return
true if decoding fails; if it succeeds, return the negation of the circuit's zero-test
on $r$ when the decoded circuit is well formed, and return true otherwise.
\end{definition}

\begin{theorem}[Testing length under size bounds]
\lean{Randomized.testLen_le_polyLen_of_le}
\label{Randomized.testLen_le_polyLen_of_le}
\leanok
\uses{Randomized.inputBits, Randomized.modParam, Randomized.numTrials, Randomized.polyLen, Randomized.testLen}
\difficulty{2}
Let $n,m,N$ be natural numbers with $n\le N$ and $m\le N$. Then $n(m+4)+16(m+12)^2\le
2400(N+1)^2$.
\end{theorem}

\begin{theorem}[Polynomial bound on testing length]
\lean{Randomized.testLen_le_polyLen}
\label{Randomized.testLen_le_polyLen}
\leanok
\uses{Randomized.AlgCircuit, Randomized.AlgCircuit.encode, Randomized.AlgCircuit.lt_length_encode, Randomized.AlgCircuit.size, Randomized.polyLen, Randomized.testLen, Randomized.testLen_le_polyLen_of_le}
\difficulty{4}
Let $c$ be a finite algebraic circuit with $n$ input variables, encoded as a binary word
of length $L$, and let $m$ be a natural number. Then $n(m+4)+16(m+12)^2\le 2400(L+1)^2$.
\end{theorem}

\begin{theorem}[Zero polynomial verifier specification]
\lean{Randomized.zeropVerifier_spec}
\label{Randomized.zeropVerifier_spec}
\leanok
\uses{Randomized.AlgCircuit.WellFormed, Randomized.AlgCircuit.decode, Randomized.AlgCircuit.encode_of_decode, Randomized.AlgCircuit.toMvPolynomial, Randomized.AlgCircuit.zeroTest, Randomized.AlgCircuit.zeroTest_of_toMvPolynomial_eq_zero, Randomized.AlgCircuit.zeroTest_sound, Randomized.ZEROP, Randomized.mem_ZEROP_iff, Randomized.polyLen, Randomized.randProb, Randomized.randProb_congr, Randomized.randProb_true, Randomized.testLen_le_polyLen, Randomized.zeropVerifier}
\difficulty{5}
Let $x$ be a finite Boolean string, and let $L$ be the language of Boolean strings
encoding well-formed arithmetic circuits over the integers whose output polynomial is
identically zero. On a uniformly random Boolean string $r$ of length $2400(|x|+1)^2$,
the verifier accepts with probability at least $2/3$ if $x\notin L$; if $x\in L$, it
rejects for every such $r$.
\end{theorem}

\begin{theorem}[Zero-polynomial testing in coRP]
\lean{Randomized.inCoRP_ZEROP}
\label{Randomized.inCoRP_ZEROP}
\leanok
\uses{Randomized.InCoRP, Randomized.VerifierModel, Randomized.ZEROP, Randomized.boolVerifier, Randomized.zeropVerifier, Randomized.zeropVerifier_spec}
\difficulty{3}
Let $E$ be a verifier model, and suppose its efficiency predicate declares the specified
zero-testing verifier efficient. This verifier accepts malformed or ill-formed circuit
encodings, and for a well-formed circuit accepts exactly when the circuit fails its
modular zero tests. The language of encodings of nonempty, well-formed arithmetic
circuits over the integers that compute the zero polynomial belongs to coRP relative to
$E$: a verifier for its complement accepts with probability at least $2/3$ on every
input in the complement and rejects every input in the language with certainty.
\end{theorem}
